%----------------------------------------------
%                中英文封面信息
%----------------------------------------------

% （1）论文信息

% 定义申请学位
% 请注意：博士学位候选人输入 Doctor，硕士学位候选人输入 Master
\def\degree{Master} 

% 论文中文标题
% 请注意：标题控制在36个汉字以内），其中用 \\ 实现标题换行，单行不得超过18个汉字。
% \underline命令用于实现标题的下划线，请把标题输入在\underline命令的的{}中
\cTitle{\underline{西南交通大学硕（博）士研究生}\\ \underline{学位论文~\LaTeX{}~模板}}

% 论文英文标题（72个字符以内，包含空格）
\eTitle{\LaTeX{}~TEMPLATE OF SOUTHWEST JIAOTONG UNIVERSITY GRADUATE STUDENT THESIS}

% 国内图书分类号
\CI{TM30}
% 国际图书分类号
\UDC{621.3}
% 保密等级
\secLevel{公开}

% （2）学生及导师信息

% 中英文姓名
\author{姜维} % 盲审使用:\author{\#\#\#}
\eAuthor{Wei Jiang} % 盲审使用:\eAuthor{\#\#\#}

% 年级
\grade{二〇二一级}
\gradeN{2021}

% 学科（工学、理学、社会学，etc）
\cDiscipline{法学}
\eDiscipline{Law}

% 专业（10个汉字以内）
% 请注意：对于超过10个汉字的专业，比如“防灾减灾工程及防护工程”，目前排版仍然不够美观，正在修正。也请相应专业的同学提供过往师兄的硕士论文以供参考。
\cMajor{马克思主义基本原理}
\eMajor{Fundamental Principle of Marxism}

% 导师的中英文姓名
% 请注意：如果需要加上导师职称，则导师姓名和导师职称之间加上\hspace{0.1em}以增大间距，整体显得更为美观，例如：诸葛亮 \hspace{0.1em} 副教授。 
\cTutor{诸葛亮} % 盲审使用:\cTutor{***}
\eTutor{Geliang Zhu} % 盲审使用:\eTutor{***}

% 封面日期，依次为年、月、日
\cDate{二零二四}{五}{二十八}
\eDate{2024}{May}{28}

%-------------------------------------------------
%                判断学位论文类型（硕士学位or博士学位）
%-------------------------------------------------

% 区分硕博类型（模板自动区分，在以上的“中英文封面信息”中进行设置）
% 请注意：本文件已经定义好需要使用的字符串，不建议用户更改。
\newif\ifdegreedoctor
\newif\ifdegreemaster

\def\temp{Doctor}
\ifx\temp\degree
\degreedoctortrue  \degreemasterfalse
\fi

\def\temp{Master}
\ifx\temp\degree
\degreedoctorfalse  \degreemastertrue
\fi

\ifdegreedoctor
\newcommand{\fitDegree}{博\quad 士}
\newcommand{\cDegree}{博士}
\newcommand{\eDegree}{Doctor of }
\fi

\ifdegreemaster
\newcommand{\fitDegree}{硕士}
\newcommand{\cDegree}{硕士}
\newcommand{\eDegree}{Master of }
\fi